\documentclass[a4paper]{article}
\usepackage{amsfonts}
\usepackage{amsmath}
\usepackage{amssymb}
\usepackage{graphicx}     

\begin{document}
  
\noindent \large Auteurs: CouldBeMathijs\\
\noindent \large Studierichting: Informatica\\
\noindent \large Oefeningenreeks 2E nr. 26.38\\

\medskip

\normalsize

$4^{\log_x 3} - 4 = 16^{\log_x \sqrt[4]{3}} + 2^{\log_x(3x)}$\\

$\Leftrightarrow 2^{2\log_x 3} - 4 = 2^{4\log_x \sqrt[4]{3}} + 2^{\log_x 3 + \log_x x}$\\

$\Leftrightarrow 2^{2\log_x 3} - 4 = 2^{2\log_x 3} + 2 \cdot 2^{2\log_x 3}$\\

\textbf{Neem} $y = 2^{2\log_x 3}$\\

$\Leftrightarrow y^2 - 4 = y + 2y$\\

$\Leftrightarrow y^2 -3y - 4 = 0$\\

$D = b^2 - 4ac = 25$\\

$\Leftrightarrow y = \dfrac{3\pm 5}{2}$\\

$\Leftrightarrow y = 4$ of $y = -1$\\

Opnieuw invullen:\\

$-1 = 2^{2\log_x 3}$: Geen oplossingen in $\mathbb{R}$\\

of\\

$4 = 2^{2\log_x 3}$\\

$\Leftrightarrow 2^2 = 2^{2\log_x 3}$\\

$\Leftrightarrow 2 = 2\log_x 3$\\

$\Leftrightarrow x = \sqrt{3}$


\begin{thebibliography}{999}
%\bibitem{a} Referentie
\end{thebibliography}
\end{document} 
